\documentclass[10pt,a4paper]{amsart}
\usepackage[margin=19mm]{geometry}
\usepackage{amsmath,amssymb,mathtools}
\usepackage[scr=rsfs,cal=euler]{mathalfa}
\usepackage{xfrac}
\usepackage[unicode,hidelinks,bookmarks=false]{hyperref}
\newcommand{\ie}{i.e.\ }
\newcommand{\cf}{cf.\ }
\newcommand{\resp}{respectively\ }
\newcommand{\Subsection}{\subsection}
\providecommand{\nobreakdash}{\nobreak\hbox{-}}
\setlength{\emergencystretch}{2em}
\multlinegap0pt
\usepackage{mathtools}
\usepackage{amssymb}
\usepackage{bm}
\usepackage{float}
\usepackage[shortlabels]{enumitem}
\usepackage{tikz}
\usepackage{xcolor}
\newtheorem{lemma}{Lemma}
\usepackage{cleveref}
\crefname{lemma}{Lemma}{Lemmas}
\newcommand{\R}{\mathbb{R}}
\DeclareMathOperator{\Ric}{Ric}
\newcommand{\bangle}[1]{\langle #1 \rangle}
\newcommand{\dvol}{\,\mathrm{d}\mu}%\mathit{vol}}
\title{Minimal hypersurfaces of Morse index one}
\author{Otis Chodosh, Matilde Gianocca}
\date{Source: arXiv:2607.27444v1 (CC BY 4.0)}
\begin{document}
\maketitle

\noindent\emph{Source and licence.} Minimal hypersurfaces of Morse index one. Version arXiv:2607.27444v1 (CC BY 4.0), distributed by arXiv under the Creative Commons Attribution 4.0 International licence, http://creativecommons.org/licenses/by/4.0/, as recorded in the arXiv licence metadata for this exact version and shown on the abstract page https://arxiv.org/abs/2607.27444v1. Full licence notice: \texttt{licenses/arxiv-2607.27444-CC-BY-4.0.txt}.\par\smallskip

\noindent\textit{Editorial scope.} Counted unit: the article's lemma lemm:sumQua-terms (source0.tex lines 651--657). The article's complete original proof of it is delivered with it (source0.tex lines 658--703, one \texttt{proof} environment of 46 lines). No mathematics is written, repaired or re-derived: the statement and the proof are the article's own lines, in the article's own notation, and no retained line is substituted or re-flowed.  Retained uncounted as the source-local context the proof needs: lines 352--354; lines 529--531; lines 533--535.  Nothing was omitted from any retained span, so the package carries no omission record.  No retained line contains a citation, so the wrapper carries no bibliography and no bibliography entry.  No retained line contains a figure environment, an image-inclusion command or drawing code, so no image is delivered and the package declares no asset dependency.  Editorial formatting only, in addition: an AMS \texttt{amsart} preamble that restates the article's own declarations and macros used by the retained lines, namely the environments \texttt{lemma} on separate counters, together with the standard packages those lines need.  The retained environments print in this document as Lemma 1 in order of appearance, whereas the article prints the same items with the numbering of its own full document.  The complete native source of the article as distributed in the arXiv e-print for this exact version is bundled unchanged as \texttt{sources/206348/source0.tex}.
\begin{equation}\label{eq:D-nabla}
D_X Y = \nabla_X Y + A(X,Y) \nu .
\end{equation}

\begin{equation}\label{eq:K_omega-pairing}
\bangle{K_\Omega \eta,\zeta} = \star (\Omega \wedge \eta\wedge \zeta)
\end{equation}

\begin{lemma}\label{lemm:P-decomp}
For $\eta \in \Lambda^2\R^N$ decomposable we have $\bangle{P_\Omega \eta,\eta} = |\eta|^2$. 
\end{lemma}

\begin{lemma}\label{lemm:sumQua-terms}
Assume that $\omega$ is harmonic on $M$. Then
\[
\sum_{a=1}^r \left(|\nabla u^\omega_a|^2 - |A|^2 (u_a^\omega)^2\right) \dvol = \frac 12 \Delta |\omega|^2 \dvol + 2 (-1)^n d( \Omega_M \wedge \omega \wedge A \omega)
\]
along $M$. 
\end{lemma}

\begin{proof}
We will work at $p\in M$. Fix $X,Y \in T_pM$ arbitrary and $e_1,\dots,e_n\in T_pM$ so that $e_1,\dots,e_n,\nu$ is an oriented orthonormal frame of $\R^N$. We compute 
\begin{equation}\label{eq:dAw-codazzi}
\begin{split}
d(A\omega)(X,Y) & = \nabla_{X}(A\omega)(Y) - \nabla_{Y}(A\omega)(X)\\
& = \omega((\nabla_{X}A) Y - (\nabla_{Y} A)X) + (\nabla_{X}\omega)(AY) - (\nabla_{Y}\omega)(AX)\\
& = (\nabla_{X}\omega)(AY) - (\nabla_{Y}\omega)(AX)\\
& = (\nabla_{AY}\omega)(X) - (\nabla_{AX}\omega)(Y)\\
& = \sum_{i=1}^n \left( \nabla_{e_i} \omega(X) \bangle{AY,e_i} - \nabla_{e_i} \omega(Y) \bangle{AX,e_i} \right)\\
& = \sum_{i=1}^n \left( \nabla_{e_i} \omega(X) \bangle{Ae_i,Y} - \nabla_{e_i} \omega(Y) \bangle{Ae_i,X}\right)\\
& = \left( \sum_{i=1}^n (\nabla_{e_i}\omega) \wedge (Ae_i)^\flat \right) (X,Y).
\end{split}
\end{equation}
We used the Codazzi equations in the third equality, $d\omega = 0$ in the fourth, and symmetry of the second fundamental form (shape operator) in the sixth. 

We now let $\beta = \nu \wedge \omega^\sharp$ and note that $\beta$ is decomposable at each point on $M$. Using \eqref{eq:D-nabla} we have 
\[
D_{e_j} \beta = - A e_j \wedge \omega^\sharp + \nu \wedge \nabla_{e_j} \omega^\sharp. 
\]
Note that $D_{e_j}\beta$ is a sum of two decomposable forms and that
\[
|Ae_j \wedge \omega^\sharp|^2 = |Ae_j|^2 |\omega|^2 - (\omega(Ae_j))^2.
\]
Using this, we find 
\begin{align*}
& \sum_{a=1}^r |\nabla u_a^\omega|^2 \dvol \\
& = \sum_{j=1}^n \bangle{P_\Omega D_{e_j} \beta,D_{e_j} \beta }\dvol  \\
& = \sum_{j=1}^n   |D_{e_j}\beta|^2 \dvol +  \sum_{j=1}^n\bangle{K_\Omega D_{e_j} \beta,D_{e_j} \beta}\dvol \\
& = (|\nabla \omega|^2  + |A|^2 |\omega|^2 -  |A\omega|^2) \dvol - 2 \star \left( \sum_{j=1}^n \Omega \wedge Ae_j \wedge \omega^\sharp \wedge \nu \wedge \nabla_{e_j}\omega^\sharp \right) \dvol\\
& = (|\nabla \omega|^2  + |A|^2 |\omega|^2 -  |A\omega|^2) \dvol + 2   \sum_{j=1}^n \Omega_M  \wedge \omega  \wedge  \nabla_{e_j}\omega\wedge (Ae_j)^\flat \\
& = (|\nabla \omega|^2  + |A|^2 |\omega|^2 -  |A\omega|^2) \dvol + 2    \Omega_M  \wedge \omega \wedge d(A \omega) \\
& = (|\nabla \omega|^2  + |A|^2 |\omega|^2 -  |A\omega|^2) \dvol + 2(-1)^n d(\Omega_M  \wedge \omega \wedge A \omega) ,
\end{align*}
where we used \eqref{eq:K_omega-pairing} and $\bangle{K_\Omega \eta,\eta} = 0$ for $\eta$ decomposable in the third equality, \eqref{eq:dAw-codazzi} in the second-to-last equality, and $d\Omega_M = d\omega = 0$ in the final equality. 


On the other hand, we have
\[
\sum_{a=1}^r (u_a^\omega)^2 = \bangle{P_\Omega \beta,\beta} = |\beta|^2 = |\omega|^2,
\]
using \Cref{lemm:P-decomp}. Since $\Ric = - A^2$ by the Gauss equations, the Bochner formula gives 
\[
\frac 12 \Delta |\omega|^2 = |\nabla \omega|^2 - |A\omega|^2 . 
\]
Combining these expressions with the above calculation proves the assertion. 
\end{proof}

\end{document}
